\documentclass[11pt,a4paper]{article}
\usepackage[T1]{fontenc}
\usepackage{lmodern}
\usepackage[hmargin=27mm,vmargin=25mm]{geometry}
\usepackage{amsmath,amssymb,amsthm}
\usepackage[hidelinks]{hyperref}
\hypersetup{
  pdftitle={Testing an ordinary-prime criterion at the first primes above fourth powers},
  pdfauthor={Robert Polak},
  pdfsubject={A sampled short-window criterion for the Riemann hypothesis}}
\setlength{\emergencystretch}{2em}
\newtheorem{theorem}{Theorem}[section]
\newtheorem{proposition}[theorem]{Proposition}
\newtheorem{lemma}[theorem]{Lemma}
\newcommand{\e}{\mathrm e}
\newcommand{\Rea}{\operatorname{Re}}
\newcommand{\dd}{\,\mathrm d}

\title{Testing an ordinary-prime criterion\\at the first primes above fourth powers}
\author{Ing.\ Robert Polak\thanks{Correspondence:
\href{mailto:robopol@gmail.com}{robopol@gmail.com}.}\\[0.4em]
\small Mathematical derivations and manuscript preparation with GPT ASTRA 6.2 (OpenAI)\\
\small Review of the underlying working note with GPT-6 Pro (OpenAI)}
\date{21 November 2026 \quad (research note, version 1.0)}

\begin{document}
\maketitle
\begin{abstract}
Let \(P(y)\) be the dyadically weighted sum of \(\log p\) over primes
between \(y\) and \(4y\), with weight one at \(2y\). We show that the
Riemann hypothesis is equivalent to the following sign test.
For arbitrarily large \(X\), one has \(P(p_m/4)\le (p_m/2)\log 2\)
at every sampled prime \(p_m\in[X,X\log X]\), where \(p_m\) is the
first prime at or above \(m^4\). A three-point difference of the
integrated von Mangoldt residual gives the weighted prime sum and an
explicit contribution from higher prime powers. The difference retains
every nontrivial zeta zero. A positive xi kernel detects a zero to
the right of the critical line inside each sufficiently late window,
and interpolation between consecutive samples bounds the cost of
checking the sign only on the prescribed sequence.
\end{abstract}

\section{The sign test}

Throughout, \(\log\) denotes the natural logarithm. For \(v>0\) and
\(y>0\), define
\begin{equation}\label{eq:prime-sum}
 w(v)=
 \begin{cases}
  2-2/v&1<v\le2,\\
  4/v-1&2<v<4,\\
  0&\text{otherwise},
 \end{cases}
 \qquad
 P(y)=\sum_{\substack{y<p<4y\\ p\ {\rm prime}}}(\log p)\,w(p/y).
\end{equation}
The weight in \eqref{eq:prime-sum} is continuous, with \(w(1)=w(4)=0\) and \(w(2)=1\).
Here \((y,4y)\) specifies the primes entering a single sum. An interval
of test positions specifies how many of these sums must satisfy the bound.

\begin{theorem}\label{thm:samples}
For each positive integer \(m\), let \(p_m=\min\{p\ge m^4:p\text{ prime}\}\).
The Riemann hypothesis (RH) is equivalent to the existence of arbitrarily
large real \(X\) for which
\begin{equation}\label{eq:sample-sign}
 P(p_m/4)\le \frac{p_m}{2}\log 2
 \qquad\text{for every }m\text{ with }X\le p_m\le X\log X .
\end{equation}
\end{theorem}

The ordinary-prime sign test originates in the work of
\texttt{bramble\_trace}, whose dyadic formulation includes tests at
integer squares \cite{bramble}. Crossley's Lean theorem, in version
1.0.4, proves one direction on real power windows. For any single
fixed \(\delta>0\), the bound \(P(x/4)\le(x/2)\log2\) for every real
\(x\in[X,X^{1+\delta}]\), for arbitrarily large \(X\), implies RH
\cite{crossley}. Crossley pointed out the connection with the
prime-power residual in his proposed addition to the research archive
\cite{pr}. The shorter interval \([X,X\log X]\) and the sequence
\((p_m)\) in Theorem~\ref{thm:samples} require an additional argument.
We give it here using the xi-kernel method of \cite{windows}.

\section{Removing higher powers}

Let \(\Lambda(p^j)=\log p\) for \(j\ge1\), with \(\Lambda(n)=0\)
at other positive integers, and let \(\gamma_E\) be Euler's constant.
Write
\begin{align*}
 F(y)&=\sum_{n\le y}\Lambda(n)(y/n-1)-y\log y+(1+\gamma_E)y,\\
 \Delta F(y)&=F(4y)-3F(2y)+2F(y),\qquad Q(y)=P(y)-2y\log 2 .
\end{align*}

\begin{lemma}\label{lem:powers}
For \(y>0\),
\begin{equation}\label{eq:exact}
 \Delta F(y)=Q(y)+H(y),\qquad
 H(y)=\sum_{j\ge2}\sum_p(\log p)\,w(p^j/y)\ge0 .
\end{equation}
As \(y\to\infty\), \(H(y)=(6\sqrt2-8+o(1))\sqrt y\).
\end{lemma}
\begin{proof}
For each \(n\), the difference applied to \((y/n-1)_+\) is \(w(n/y)\),
as follows by splitting at \(n=y,2y,4y\). The linear term cancels and
\(\Delta(-y\log y)=-2y\log2\). Separating the primes gives
\eqref{eq:exact}. At \(y=2\), for instance, the square \(2^2\) has
weight one and accounts for all of \(H(2)=\log2\).

For squares, the prime number theorem for
\(\vartheta(v)=\sum_{p\le v}\log p\) and partial summation give
\[
 \frac{1}{\sqrt y}\sum_p(\log p)\,w(p^2/y)
 \ \longrightarrow\ \int_1^2 w(u^2)\dd u
 =\int_1^{\sqrt2}(2-2/u^2)\dd u+
   \int_{\sqrt2}^{2}(4/u^2-1)\dd u
 =6\sqrt2-8 .
\]
Indeed \(\vartheta(u\sqrt y)/\sqrt y\to u\) uniformly for \(1\le u\le2\),
and \(u\mapsto w(u^2)\) has bounded variation. An exponent \(j\ge3\)
contributes at most \((4y)^{1/3}\log(4y)\), by counting possible bases
and using \(0\le w\le1\). There are \(O(\log(4y))\) such exponents.
Their total is \(o(\sqrt y)\).
\end{proof}

\section{The effect of the difference on zeta zeros}

Use the spectral sum from \cite{cone,windows},
\[
 B(t)=\sum_\rho a_\rho\e^{(\rho-1/2)t},\qquad
 a_\rho=\frac1{\rho(1-\rho)},\qquad
 C=B(0)=2+\gamma_E-\log(4\pi).
\]
Here the nontrivial zeros are counted with multiplicity and
\(\sum_\rho |a_\rho|<\infty\). Put \(\ell=\log2\) and
\[
 B_\Delta(t)=2B(t+2\ell)-3\sqrt2\,B(t+\ell)+2B(t).
\]
The explicit formula in \cite[Section~2]{cone} gives, for \(y>1\),
\[
 F(y)=\sqrt y\,B(\log y)+\log(2\pi)+R(y),\qquad
 R(y)=-\sum_{j\ge1}\frac{y^{-2j}}{2j(2j+1)}
     =-\frac1{6y^2}+O(y^{-4}).
\]
Since \(1-3+2=0\), the constant \(\log(2\pi)\) cancels.
The leading coefficient of \(\Delta R\) is
\(-(4^{-2}-3\cdot2^{-2}+2)/6=-7/32\). Thus
\begin{equation}\label{eq:remainder}
 \frac{\Delta F(y)}{\sqrt y}
 =B_\Delta(\log y)-\frac7{32}y^{-5/2}+O(y^{-9/2}).
\end{equation}
In the series for \(B_\Delta\), the coefficient of each zero is
\(a_\rho M(\rho)\), with
\begin{equation}\label{eq:multiplier}
 M(z)=4^z-3\cdot2^z+2=(2^z-1)(2^z-2).
\end{equation}
In \(0<\Rea z<1\) we have \(1<|2^z|<2\). Hence \(M\) has no zero
in the open strip, and \(|M(z)|\le12\) there. Crossley's delayed
transform uses the symbol \(-4^{-z}M(z)\), corresponding to the
opposite sign and the variable \(x=4y\) \cite{crossley}.

\begin{lemma}\label{lem:osc}
Fix a zero \(\rho_*\) with \(\alpha=\Rea\rho_*-1/2>0\).
There is \(b_*>0\) such that for every sufficiently large \(T\)
some \(t\) satisfies
\begin{equation}\label{eq:osc}
 |t-T|\le R_T=\tfrac12\log(T/2),\qquad B_\Delta(t)\ge b_*\e^{\alpha T}.
\end{equation}
\end{lemma}
\begin{proof}
We recall the ingredients of \cite[Lemma~2.1]{windows} to justify its
use after taking the difference. The positive, integrable kernel \(K\)
has bilateral transform \(\xi(1/2+z)\). For a zero \(\rho_*\) of
multiplicity \(r\) the isolating kernel \(h_r\) satisfies
\(|h_r|\le LK\) for a constant \(L>0\), and its transform is
\(\xi(1/2+z)/(z-\lambda_*)^r\), where \(\lambda_*=\rho_*-1/2\).
This transform has a nonzero value at \(\lambda_*\) and vanishes at
every other shifted zero.

The coefficients modified by \eqref{eq:multiplier} are absolutely summable. Conjugate pairing
makes \(B_\Delta\) real, and \(|B_\Delta(v)|\le C_0\e^{|v|/2}\)
on the real line for some \(C_0>0\). Termwise integration therefore
gives
\[
 \int_{\mathbb R} K(u)B_\Delta(T+u)\dd u=0,\qquad
 \int_{\mathbb R} h_r(u)B_\Delta(T+u)\dd u
    =M(\rho_*)A_*\e^{\lambda_*T},\qquad A_*\ne0 .
\]
In the second equality the \(r\) contributions from \(\rho_*\) add.
Since \(K\) is positive, the integrals of the positive and negative
parts against \(K\) are equal. Each is at least
\(|M(\rho_*)A_*|\e^{\alpha T}/(2L)\).
The bound \(K(u)\le C_1\exp(9|u|/2-\pi\e^{2|u|})\) gives
\[
 \int_{|u|>R_T}K(u)|B_\Delta(T+u)|\dd u
 =O\!\left(T^{3/2}\e^{-(\pi-1)T/2}\right)=o(1)
\]
by the same integral estimate as in \cite{windows}. Dividing the
remaining positive mass by \(\int_{\mathbb R}K(u)\dd u\) proves
\eqref{eq:osc}.
\end{proof}

\section{Short windows and interpolation}

\begin{proposition}\label{prop:windows}
Fix \(0\le s<1/2\). Suppose there is a constant \(A\ge0\) such that
for arbitrarily large \(Z\) one has
\[
 Q(y)\le A y^{1/2+s}\qquad(Z\le y\le Z\log Z).
\]
Then every nontrivial zeta zero has real part at most \(1/2+s\).
\end{proposition}
\begin{proof}
By \eqref{eq:exact}, \eqref{eq:remainder} and \(H(y)=O(\sqrt y)\),
there is a constant \(C_s>0\) such that on each such window,
for \(Z\) sufficiently large, \(B_\Delta(t)\le C_s\e^{st}\)
at every \(t=\log y\). Suppose a zero has \(\alpha>s\).
Choose \(U=\log Z\) and \(T=U+\frac12\log U\). Because \(T\le2U\)
eventually, \([T-R_T,T+R_T]\subseteq[U,U+\log U]\).
At the point from Lemma~\ref{lem:osc} we would have
\[
 b_*\e^{\alpha T}\le C_s\e^{sT}(T/2)^{s/2}.
\]
This fails as the window starts tend to infinity.
\end{proof}

\begin{lemma}\label{lem:forward}
Assuming RH, \(Q(y)<0\) for all sufficiently large real \(y\).
\end{lemma}
\begin{proof}
Under RH all \(a_\rho\) are positive and sum to \(C\).
On the critical line \(|M(\rho)|\le4+3\sqrt2\). Consequently
\(|B_\Delta(t)|\le(4+3\sqrt2)C\) uniformly. Lemma~\ref{lem:powers}
and \eqref{eq:remainder} give
\[
 \limsup_{y\to\infty}\frac{Q(y)}{\sqrt y}
 \le (4+3\sqrt2)C-(6\sqrt2-8)<-0.1045 .
\]
The strict inequality follows by substituting the displayed value of \(C\).
\end{proof}

In particular the sign on complete real windows \([X,X\log X]\) in
the variable \(x=4y\) is equivalent to RH. The forward implication
uses Lemma~\ref{lem:forward}. Conversely, taking \(Z=X/4\) gives the
sign on \([Z,Z\log Z]\), so Proposition~\ref{prop:windows} with \(s=0\)
and the functional equation give RH.

\begin{proof}[Proof of Theorem~\ref{thm:samples}]
Lemma~\ref{lem:forward} proves necessity. For sufficiency put
\(a_m=p_m/4\). The prime gap estimate of Baker, Harman and Pintz
\cite{bhp} gives \(p_m=m^4+O(m^{2.1})\) as in \cite{cone}.
Hence the sequence is eventually strictly increasing and its gaps
are \(O(a_m^{3/4})\). We first estimate the interpolation error
between two consecutive sample points \(a,b\), with \(h=b-a\).

As a function of \(y\), \(Q(y)\) is continuous and piecewise affine.
For a prime \(p\) the three slope jumps occur at \(p/4,p/2,p\) and
are \(4,-6,2\), respectively, in units of \((\log p)/p\).
Within \([a,b]\) the corresponding primes belong to
\([4a,4b]\), \([2a,2b]\), or \([a,b]\). Each interval contains
\(O(h+1)\) integers, each with \((\log p)/p=O((\log a)/a)\).
The sum of the absolute jumps is thus \(O((h+1)(\log a)/a)\).
The chord slope is an average of the one-sided slopes on the open
pieces. Integrating their difference from this average bounds the
absolute error from the chord by
\[
 O\!\left(h(h+1)\frac{\log a}{a}\right)=O(\sqrt a\log a).
\]
When \(Q(a),Q(b)\le0\), the chord is nonpositive, so
\(Q(y)\le C_2\sqrt y\log y\) for \(a\le y\le b\), with \(C_2\)
independent of the samples.

In a window satisfying \eqref{eq:sample-sign}, the admissible \(a_m\)
lie in \([Z,Z\log(4Z)]\), where \(Z=X/4\). The distances from the
ends to the nearest samples inside the window are respectively
\(O(Z^{3/4})\) and \(O((Z\log(4Z))^{3/4})\), by the gap bound.
For \(U=\log Z\) and \(T=U+\frac12\log U\) one has
\[
 \frac{\e^{T-R_T}}{Z}\longrightarrow\sqrt2,\qquad
 \frac{\e^{T+R_T}}{Z\log Z}\longrightarrow\frac1{\sqrt2}.
\]
Thus \([\e^{T-R_T},\e^{T+R_T}]\) lies between the first and last
sample of the window for large \(Z\). Interpolation, \eqref{eq:exact}
and \eqref{eq:remainder} give \(B_\Delta(t)\le C_3(1+t)\) throughout
\([T-R_T,T+R_T]\). A fixed zero with \(\alpha>0\) would give
\(b_*\e^{\alpha T}\le C_3(1+T+R_T)\) on arbitrarily late windows.
This contradiction excludes every zero to the right of the line,
and the functional equation proves RH.
\end{proof}

The arithmetic question is whether windows satisfying
\eqref{eq:sample-sign} can be constructed directly from the primes.
Theorem~\ref{thm:samples} establishes the equivalence. The
shorter-window and sampled statements in this note are written
mathematical proofs. Crossley's Lean declaration cited above
formalises the real power-window implication.

\begin{thebibliography}{9}
\bibitem{bramble} \texttt{bramble\_trace}, \emph{Reading Alp\"oge's
proposed proof of RH from theta lifts to prime sums}, rStation, 2026.
\href{https://rstation.ai/r/riemann_hypothesis/posts/100}{Author's discussion}
and \href{https://rstation.ai/api/v1/posts/uw5cfb_J_Cj7XJO6/attachment?revision=24}{revision 24}.
\bibitem{crossley} L.~Crossley, \emph{An ordinary-prime window criterion
for the Riemann hypothesis}, version 1.0.4, 2026.
\href{https://doi.org/10.5281/zenodo.24412819}{Zenodo}.
\href{https://github.com/leecrossley/prime-window-rh-criterion/blob/v1.0.4/PrimeWindow.lean#L195-L211}{Lean declaration}.
\bibitem{pr} L.~Crossley, \emph{Relate the short-window bound to the
ordinary-prime test}, 2026.
\href{https://github.com/robopol/Riemann-hypothesis/pull/3}{RoboPol research archive, PR 3}.
\bibitem{cone} R.~Polak, \emph{An Explicit Prime-Power Cone Criterion for
the Riemann Hypothesis}, version 1.2, 2026.
\href{https://doi.org/10.5281/zenodo.24406562}{Zenodo}.
\bibitem{windows} R.~Polak, \emph{Importing arithmetic windows into the
prime-power cone}, version 1.0, 2026.
\href{https://doi.org/10.5281/zenodo.24415724}{Zenodo}.
\bibitem{bhp} R.~C.~Baker, G.~Harman and J.~Pintz,
The difference between consecutive primes, II,
\emph{Proc. London Math. Soc.} (3) \textbf{83} (2001), 532--562.
\href{https://doi.org/10.1112/plms/83.3.532}{DOI}.
\end{thebibliography}
\end{document}
